\documentclass[10pt,a4paper]{amsart}
\usepackage[margin=19mm]{geometry}
\usepackage{amsmath,amssymb,mathtools}
\usepackage[scr=rsfs,cal=euler]{mathalfa}
\usepackage{xfrac}
\usepackage[unicode,hidelinks,bookmarks=false]{hyperref}
\newcommand{\ie}{i.e.\ }
\newcommand{\cf}{cf.\ }
\newcommand{\resp}{respectively\ }
\newcommand{\Subsection}{\subsection}
\providecommand{\nobreakdash}{\nobreak\hbox{-}}
\setlength{\emergencystretch}{2em}
\multlinegap0pt
\usepackage{amssymb}
\usepackage{tcolorbox}
\usepackage{mathtools}
\usepackage{float}
\usepackage{xcolor}
\usepackage{bm}
\newtheorem{proposition}{Proposition}
\newtheorem{lemma}{Lemma}
\title{\bf superintegrable systems from Cartan commutants}
\author{Kai Jiang, Guorui Ma, Ian Marquette, Junze Zhang, Yao-Zhong Zhang}
\date{Source: arXiv:2601.01369v2 (CC BY 4.0)}
\begin{document}
\maketitle

\noindent\emph{Source and licence.} \bf superintegrable systems from Cartan commutants. Version arXiv:2601.01369v2 (CC BY 4.0), distributed by arXiv under the Creative Commons Attribution 4.0 International licence, http://creativecommons.org/licenses/by/4.0/, as recorded in the arXiv licence metadata for this exact version and shown on the abstract page https://arxiv.org/abs/2601.01369v2. Full licence notice: \texttt{licenses/arxiv-2601.01369-CC-BY-4.0.txt}.\par\smallskip

\noindent\textit{Editorial scope.} Counted unit: the article's lemma alg1 (source0.tex lines 1646--1661). The article's complete original proof of it is delivered with it (source0.tex lines 1663--1717, one \texttt{proof} environment of 55 lines). No mathematics is written, repaired or re-derived: the statement and the proof are the article's own lines, in the article's own notation, and no retained line is substituted or re-flowed.  Retained uncounted as the source-local context the proof needs: lines 850--861; lines 1611--1621.  Nothing was omitted from any retained span, so the package carries no omission record.  No retained line contains a citation, so the wrapper carries no bibliography and no bibliography entry.  No retained line contains a figure environment, an image-inclusion command or drawing code, so no image is delivered and the package declares no asset dependency.  Editorial formatting only, in addition: an AMS \texttt{amsart} preamble that restates the article's own declarations and macros used by the retained lines, namely the environments \texttt{proposition}, \texttt{lemma} on separate counters, together with the standard packages those lines need.  The retained environments print in this document as Proposition 1, Lemma 1 in order of appearance, whereas the article prints the same items with the numbering of its own full document.  The complete native source of the article as distributed in the arXiv e-print for this exact version is bundled unchanged as \texttt{sources/192074/source0.tex}.
 \begin{proposition}
 \label{prop:slicepoisson}
Let $h_1,h_2 \in S(\mathfrak{g})$ and $\theta_1,\theta_2\in S(\mathfrak{m}-\varepsilon W)^A$. Then %$ \in P^*(S(\mathfrak{g}))$ $\,\{\pi_\mathfrak{m}^*\theta_1,\pi_\mathfrak{m}^*\theta_2\}_\varepsilon\in \pi_\mathfrak{m}^* S(\mathfrak{m}-\varepsilon W)^A$ and 
\begin{align*}
\{P^*h_1,P^*h_2\}_\varepsilon = P^*\left(\{h_1,h_2\}_1\right) \text{ and } \{\pi_\mathfrak{m}^*\theta_1,\pi_\mathfrak{m}^*\theta_2\}_\varepsilon=\pi_\mathfrak{m}^*\bigl(\{\theta_1,\theta_2\}_2\bigr),
\end{align*} where $\{\cdot,\cdot\}_1 = \{\cdot,\cdot\} $ is a Lie-Poisson bracket and $\{\cdot,\cdot\}_2$ is the Poisson bracket on $S(\mathfrak{m}-\varepsilon W)^A$ given by \begin{align*}
\{\theta_1,\theta_2\}_2(\xi) :=-B \left(\xi,\bigl[(\nabla\theta_1(\xi))_{\mathfrak{m}},(\nabla\theta_2(\xi))_{\mathfrak{m}}\bigr]\right),
\qquad \xi\in \mathfrak{m}-\varepsilon W .
\end{align*}
In particular, $\mathfrak{F}_2^{\mathrm{poly}}=\pi_\mathfrak{m}^* S(\mathfrak{m}-\varepsilon W)^A$
is a Poisson subalgebra of $\mathcal{O}(T^* M)$.
\end{proposition}

\begin{lemma} 
\label{lem:root-coord}
Let $\mathfrak{g}_\mathbb{C}=\mathfrak{sl}_3(\mathbb{C})$ with the Cartan subalgebra $\mathfrak{t}_\mathbb{C}$ and its root space decomposition $\mathfrak{g}_\mathbb{C} = \mathfrak{t}_\mathbb{C}  \oplus \bigoplus_{\alpha\in\Delta}\mathfrak{g}_\alpha$ with $  \dim \mathfrak{g}_{\pm\alpha}=1$. Choose root vectors $E_{\pm\alpha_k}\in\mathfrak{g}_{\pm\alpha_k}$ for the three positive roots $\alpha_1,\alpha_2,\alpha_3=\alpha_1+\alpha_2$, and the Killing form $B$ such that
\begin{align}
B(E_{\alpha_k},E_{-\alpha_\ell})=\delta_{k\ell},\qquad B(\mathfrak{t}_\mathbb{C},\mathfrak{g}_{\pm\alpha})=0,\qquad B(\mathfrak{g}_\alpha,\mathfrak{g}_\beta)=0 \text{ if }\alpha+\beta\neq 0. \label{eq:normal} 
\end{align} Take $E_{\pm \alpha_k}$ in $\mathfrak{sl}(3,\mathbb{C})$ such that the normalization in \eqref{eq:normal} holds. Let $\mathfrak{g}=\mathfrak{su}(3)$ be the compact real form with the conjugation $E_{-\alpha_k}=\overline{E_{\alpha_k}}$. Then every $\zeta\in\mathfrak{g}$ admits a unique expansion \begin{align*}
\zeta = H +  \sum_{k=1}^3\Big(z_k\,E_{\alpha_k}+\overline{z_k}\,E_{-\alpha_k}\Big), \qquad H\in\mathfrak{t},\ z_k\in\mathbb{C},
\end{align*} and the coefficients are: \begin{align}
  z_k:=  z_k(\zeta) = B(\zeta,E_{-\alpha_k}),\qquad \overline{z_k} : = \overline{z_k}(\zeta)= B(\zeta,E_{\alpha_k}) \quad  k=1,2,3. \label{eq:complexcoe}
\end{align}
\end{lemma}

 \begin{lemma}
 \label{alg1}
The Poisson algebra $S(\mathfrak{m})^T$ of the $T$-invariant polynomials with a Poisson structure $\{\cdot,\cdot\}_2$ induced by $\{\cdot,\cdot\}_\varepsilon$ is minimally generated by the five polynomials as follows:  \begin{align*}
    u_k=|z_k|^2\,, \text{ } k=1,2,3 ,\qquad v=\Re(z_1z_2\bar z_3),\qquad w=\mathrm{Im}(z_1z_2\bar z_3),
\end{align*} subject to the single cubic relation $ u_1u_2u_3=v^2+w^2 $.  Then   \begin{align*}
    \mathfrak{F}_2^{\mathrm{poly}} =  \pi_\mathfrak{m}^* \left(\mathbb{R}[u_1,u_2,u_3,v,w]\Big/ \langle u_1u_2u_3-v^2-w^2\rangle\right), \qquad \deg u_k=2,\,\deg v=\deg w=3.
\end{align*}  The Poisson brackets among these generators are  \begin{align}
 \nonumber
     &\{u_1,u_2\}_2=2v,\ \{u_2,u_3\}=2v,\ \ \ \{u_3,u_1\}_2=2v,\\
     \nonumber
&\{u_1,v\}_2=u_1(u_3-u_2)-c_1w,\quad    \{u_2,v\}_2=u_2(u_1-u_3)-c_2w,\\
&\{u_3,v\}_2=u_3(u_2-u_1)-c_3w,\quad   \{u_k,w\}_2=c_k v,\\
  \nonumber
&\{v,w\}_2=-\frac{1}{2}\big(c_3u_1u_2-c_2u_1u_3-c_1u_2u_3\big),
 \end{align}  where $ 1\leq  k \leq 3$ and $c_k = \varepsilon B(W,H_{\alpha_k})$. %This makes the generic Poisson rank equal to $2$.
\end{lemma}

\begin{proof}
We first compute the $T$-invariant generators in $S(\mathfrak{m})$. By Lemma \ref{lem:root-coord}, we deduce  \begin{align}
    u_k:= |z_k|^2=x_k^2+y_k^2\quad \text{ with } k=1,2,3,\qquad v+\mathrm{i}w:=z_1\,z_2\,\overline{z_3} . \label{eq:tinvari}
\end{align} We first show that the polynomials defined in \eqref{eq:tinvari} are $T$-invariant. For the degree two generators, it is obvious for $u_k$ as $\alpha_k +  (-\alpha_k) = 0$. Moreover, for $z_1 z_2 \overline{z_3}$, we use the root identity $\alpha_1+\alpha_2-\alpha_3=0$ since the total weight is zero. Hence, it is $T$-invariant. We then take its real and imaginary parts $v=\Re(\cdot)$, $w=\mathrm{Im}(\cdot)$. A general $T$-invariant monomial is $z_1^{a_1}z_2^{a_2}z_3^{a_3}
\overline{z_1}^{b_1}\overline{z_2}^{b_2}\overline{z_3}^{b_3}$ with weight $(a_1-b_1)\alpha_1+(a_2-b_2)\alpha_2+(a_3-b_3)\alpha_3$. That is, by the definition of the Cartan commutant \begin{align*}
  (a_1-b_1)\alpha_1+(a_2-b_2)\alpha_2+(a_3-b_3)(\alpha_1+\alpha_2)=0,  
\end{align*} which is equivalent to the linear system $ (a_1-b_1)+(a_3-b_3)=0$ and  $ (a_2-b_2)+(a_3-b_3)=0$. Hence, $a_1-b_1=-(a_3-b_3)$ and $a_2-b_2=-(a_3-b_3)$. Writing $d:=a_3-b_3$, we can decompose any invariant monomial as a product of the powers of $|z_k|^2=u_k$ (which accounts for the symmetric part) and $(z_1z_2\overline{z_3})^d$ (which accounts for the weight $d$), together with its complex conjugate, if needed, to keep the coefficients real. Thus, the invariant ring is generated by $u_1,u_2,u_3$ and $z_1z_2\overline{z_3}$ (hence, by $u_1,u_2,u_3,v,w$).  In addition, a routine computation shows that there is an algebraic relation between these five invariants: \begin{align}
 u_1\,u_2\,u_3=|z_1|^2\,|z_2|^2\,|z_3|^2 =|z_1 z_2 \overline{z_3}|^2=v^2+w^2 . \label{eq:plrelation}
\end{align}   % The single algebraic relation is the one that we have shown in \eqref{eq:plrelation}, with both sides equal to $|z_1 z_2 \overline{z_3}|^2$.

 
   We then calculate the Poisson bracket relations from all the $5$ generators $\{u_1,u_2,u_3,w,v\}$. By Proposition \ref{prop:slicepoisson}, the Poisson bracket defined on $S(\mathfrak{m})^T$ is given by $\{\theta_1,\theta_2\}_2(\xi) :=-B \left(\xi,\bigl[(\nabla\theta_1(\xi))_{\mathfrak{m}},(\nabla\theta_2(\xi))_{\mathfrak{m}}\bigr]\right)$, where $(\cdot)_\mathfrak{m}$ is projection to $\mathfrak{m}$-component. In other words, we need to evaluate the gradient of the coordinate functions. By Lemma \ref{lem:root-coord}, the linear functional of coordinates is given by $z_k(\zeta)=B(\zeta,E_{-\alpha_k})$, and $\overline{z_k}(\zeta)=B(\zeta,E_{\alpha_k})$ are linear functionals. A direct calculation shows that  $\nabla z_k=E_{-\alpha_k}$ and $\nabla \overline{z_k}=E_{\alpha_k}$. Therefore, \begin{align}
\nabla u_k=\nabla(z_k\overline{z_k})=\overline{z_k}\,E_{\alpha_k}+z_k\,E_{-\alpha_k}.\tag{A}
\end{align} Writing $F:=z_1z_2\overline{z_3}$ such that $v=\frac{1}{2}(F+\overline{F})$, $w=\frac{1}{2\mathrm{i}}(F-\overline{F})$, \begin{align}
\nabla F=z_2\overline{z_3}\,E_{-\alpha_1}+z_1\overline{z_3}\,E_{-\alpha_2}+z_1z_2\,E_{\alpha_3},\qquad \nabla\overline{F}=\overline{z_2}z_3\,E_{\alpha_1}+\overline{z_1}z_3\,E_{\alpha_2}+\overline{z_1}\,\overline{z_2}\,E_{-\alpha_3}.\tag{B}
\end{align} From (A) and (B), we first deduce $(\theta)_\mathfrak{m} =\theta$ for any $\theta \in S(\mathfrak{m})^T$.

We now compute the Poisson bracket $\{u_1,u_2\}_2$.  Using (A) and (B), \begin{align*}
[\nabla u_1,\nabla u_2] = z_1z_2\,[E_{-\alpha_1},E_{-\alpha_2}]+\overline{z_1}\,\overline{z_2}\,[E_{\alpha_1},E_{\alpha_2}] =  \overline{z_1}\,\overline{z_2}\,E_{\alpha_3}-\,z_1z_2\,E_{-\alpha_3}.
\end{align*} Therefore, by Lemma \ref{lem:root-coord} and its remark, we have $B(\zeta,E_{\alpha_3}) = -z_3$ and $B(\zeta,E_{-\alpha_3}) = \overline{z_3}$, and hence \begin{align*}
\{u_1,u_2\}_2 =-B\big(\zeta,\,[\nabla u_1,\nabla u_2]\big) = -\Big( \overline{z_1}\,\overline{z_2}\,B(\zeta,E_{\alpha_3}) -z_1z_2\,B(\zeta,E_{-\alpha_3})\Big) = \overline{z_1}\,\overline{z_2}\,z_3 +z_1z_2\,\overline{z_3} = 2\,\Re(z_1z_2\overline{z_3}) = 2\,v.
\end{align*}

 Now, we focus on the Poisson bracket $\{u_1,F\}_2$. Similar to what we did above, a direct computation shows that \begin{align*}
[\nabla u_1,\nabla F] &= \left[\,\overline{z_1}E_{\alpha_1}+z_1E_{-\alpha_1},\ z_2\overline{z_3}E_{-\alpha_1}+z_1\overline{z_3}E_{-\alpha_2}+z_1z_2E_{\alpha_3}\right]\\
&= \overline{z_1}\,z_1z_2\,[E_{\alpha_1},E_{\alpha_3}]+\overline{z_1}\,z_1\overline{z_3}\,[E_{\alpha_1},E_{-\alpha_2}] + \overline{z_1}\,z_2\overline{z_3}\,[E_{\alpha_1},E_{-\alpha_1}]\\
&\quad +  z_1\,z_1\overline{z_3}\,[E_{-\alpha_1},E_{-\alpha_2}] +  z_1\,z_1z_2\,[E_{-\alpha_1},E_{\alpha_3}] +  z_1\,z_2\overline{z_3}\,[E_{-\alpha_1},E_{-\alpha_1}]\\
&=  \underbrace{z_2\overline{z_3}\,\overline{z_1}}_{\;=\ \overline{z_1}z_2\overline{z_3}}\,H_{\alpha_1} - \underbrace{z_1\overline{z_1}\,\overline{z_3}}_{\;=\ |z_1|^2\overline{z_3}}\,E_{-\alpha_3} + \underbrace{z_1\overline{z_1}\,z_2}_{\;=\ |z_1|^2 z_2}\,E_{\alpha_2}.
\end{align*}
Then the twisted Poisson bracket $\{u_1,F\}_2$ is calculated as follows: \begin{align*}
\{u_1,F\}_2 =-\,B(\zeta,[\nabla u_1,\nabla F]) = -\Big(\overline{z_1}z_2\overline{z_3}\,B(\zeta,H_{\alpha_1}) - |z_1|^2\overline{z_3}\,B(\zeta,E_{-\alpha_3}) +  |z_1|^2 z_2\,B(\zeta,E_{\alpha_2})\Big).
\end{align*} Using $B(\zeta,H_{\alpha_1})=B(X,H_{\alpha_1})-\varepsilon\,c_1$ and $B(\zeta,E_{-\alpha_3})=z_3,\ B(\zeta,E_{\alpha_2})=\overline{z_2}$, we obtain \begin{align}
\{u_1,F\}_2 = -\,\overline{z_1}z_2\overline{z_3}\, \left(B(X,H_{\alpha_1})-\varepsilon c_1 \right) -|z_1|^2\big( |z_2|^2-|z_3|^2 \big). \label{eq:a1}
\end{align}

 Finally, we compute the twisted Poisson bracket $\{u_1,\overline{F}\}_2$. A similar calculation from (A) and (B) gives \begin{align*}
[\nabla u_1,\nabla \overline{F}] = -\,z_1\overline{z_2}z_3\,H_{\alpha_1}+|z_1|^2 z_3\,E_{\alpha_3}-|z_1|^2\overline{z_2}\,E_{-\alpha_2}.
\end{align*} Hence, \begin{align}
\{u_1,\overline{F}\}_2 = z_1\overline{z_2}z_3\,\left(B(X,H_{\alpha_1})-\varepsilon c_1\right) -|z_1|^2\big(\,|z_3|^2-|z_2|^2\,\big). \label{eq:a2}
\end{align}

To obtain the corresponding Poisson bracket, we then need to transform from $F,\overline{F}$ to $v,w$.  Since $v=\frac{1}{2}(F+\overline{F})$ and $w=\frac{1}{2\mathrm{i}}(F-\overline{F})$, \begin{align*}
\{u_1,v\}_2 =\frac{1}{2}\big(\{u_1,F\}_2+\{u_1,\overline{F}\}_2\big),\qquad \{u_1,w\}_2 =\frac{1}{2\mathrm{i}}\big(\{u_1,F\}_2-\{u_1,\overline{F}\}_2\big).
\end{align*} Adding \eqref{eq:a1} and \eqref{eq:a2} together cancels the terms $|z_1|^2$ and yields \begin{align*}
\{u_1,v\}_2 = -\,\Re\Big(\overline{z_1}z_2\overline{z_3}-z_1\overline{z_2}z_3\Big)\left(B(X,H_{\alpha_1})-\varepsilon c_1\right) = -\,\varepsilon c_1\,\Im(z_1z_2\overline{z_3}) +|z_1|^2\big(|z_3|^2-|z_2|^2\big).
\end{align*} Recognizing $w=\Im(z_1z_2\overline{z_3})$ and $u_k=|z_k|^2$, we have \begin{align}
 \{u_1,v\}_2=u_1(u_3-u_2)-c_1\,w . \label{eq:b}
\end{align} Similarly, \begin{align}
\{u_1,w\}_2 =\frac{1}{2\mathrm{i}}\big(\{u_1,F\}_2-\{u_1,\overline{F}\}_2\big) =  \frac{c_1}{2}\big(\overline{z_1}z_2\overline{z_3}+z_1\overline{z_2}z_3\big) =  c_1 v . \label{eq:c}
\end{align} A cyclic permutation of $(1,2,3)$ yields $\{u_i,u\}_2$ and $\{u_i,w\}_2 $. Finally, \begin{align*}
\{v,w\}_2 =\frac{1}{4\mathrm{i}}\big(\{F+\overline{F},F-\overline{F}\}_2\big) =\frac{1}{2\mathrm{i}}\,\{F,\overline{F}\}_2- B \big(\zeta,\,[\nabla F,\nabla\overline{F}]\big).
\end{align*} Here, we deduce that $B\big(\zeta,\,[\nabla F,\nabla\overline{F}]\big) = \frac{\mathrm{i}}{2}\,\Big(c_3 u_1u_2-c_2 u_1u_3-c_1 u_2u_3\Big) $ such that \begin{align}
 \{v,w\}_2=-\frac{1}{2}\Big(c_3 u_1u_2-c_2 u_1u_3-c_1 u_2u_3\Big).
\end{align} This completes the proof.
\end{proof}

\end{document}
